\scanpage{21}
D'où $s_*(z^2)=t_*(\Theta^2c')$, où $t:A\times Y\unsure{-|\longrightarrow}A$ est la projection. On a alors, du pt de vue \unsure{degré}

\diagram{\begin{tikzcd}[ampersand replacement=\&] A\times Y\arrow[d,"t"']\\ A\end{tikzcd}}

\[
\begin{aligned}
\unsure{\mathrm{tr}(z^2)|\mathrm{tr}(s_*(z^2))|K(z^2)}&=\mathrm{tr}(\Theta^2c')=\mathrm{tr}(\Theta^2\cdot1_A\times c)\\
&=\mathrm{tr}(p_*(\Theta^2)\cdot c)
\end{aligned}
\]

où $p$ est la projection $A\times Y\unsure{-|\longrightarrow}Y$.

Nous allons calculer $p_*(\Theta^{\unsure{2|?}})$, en décomposant $p$ en $\deleted{\unsure{q}}p'\times\mathrm{id}_{P}$, où $p':A\times A\longrightarrow A$ est la deuxième projection. On a

\diagram{\begin{tikzcd}[ampersand replacement=\&,row sep=small]
 A\times\deleted{\unsure{X}}Y\arrow[r,"p"] \arrow[d,equal] \&Y\\
 A\times\underbrace{A\times P}\arrow[r] \&\underbrace{A\times P}
\end{tikzcd}}

\[
\begin{aligned}
\Theta&=\Theta_0\times1_P\quad\text{d'où }\Theta^2=\Theta_0^2\times1_P\qquad\text{d'où}\\
p_*(\Theta^2)&=p'_*(\Theta_0^2)\times\mathrm{id}_{P*}(1_P)=p'_*(\Theta_0^2)\times1_P.
\end{aligned}
\]

D'autre part

\[
\begin{aligned}
\Theta_0^2&=\deleted{\Delta}\,\Delta^2+0+0-2(0,0)-2(0,0)+2(0,0)\\
&=-2(0,0)
\end{aligned}
\]

compte tenu que $\Delta^2=0$ \inserted{(équiv. alg.)} puisque $\chi(A)=0$.

d'où

\[
p_*(\Theta^2)=-2(0\times0\times P)
\]

et on trouve par suite

\[
\mathrm{tr}(s_*(z^2))=-2c(0,0)
\]

où on pose

\[
c(0,0)=\mathrm{tr}\bigl(c\cdot(0\times\unsure{0|\varnothing}\times P)\bigr)
\]

Il existe un $c$ tel que $c(0,0)=1$, p. ex. $c=A\times A\times(a)$, où $a\in P(k)$. Donc remplaçant \inserted{\unsure{ensuite}} $c$ par $-c$, on trouve comme valeurs successivement $-2$ et $2$, qui ont des signes différents !!

% archive sheet number: 20
